\subsection{代数的形式光滑商}

定理 2.--- 设 k 为环，A 为 k-代数，J 为 A 的理想且 k-代数 A/J 形式光滑。赋予 A 以 J-进拓扑。下列条件等价：

\begin{enumerate}
\def\labelenumi{(\roman{enumi})}
\item
  拓扑 k-代数 A 形式光滑；
\item
  \(A/\mathrm{J}\)-模 \(\mathrm{J}/\mathrm{J}^2\) 是投射的，且典范同态 (§ 5, 第 2 节)
\end{enumerate}

\[
\beta : \mathbf {S} _ {\mathrm{A/J}} (\mathrm{J} / \mathrm{J} ^ {2}) \rightarrow \operatorname{gr} _ {\mathrm{J}} (\mathrm{A})
\]

是双射的；

\begin{enumerate}
\def\labelenumi{(\roman{enumi})}
\setcounter{enumi}{2}
\tightlist
\item
  A/J-模 J/J \(^{2}\) 是投射的，且存在从 A 的分离完备化到分次代数 \(\mathsf{S}_{A/J}\) (J/J \(^{2}\)) 的完备化上的拓扑 k-代数同构。
\end{enumerate}

若 A 是 Noether 的，这些条件也等价于：

\begin{enumerate}
\def\labelenumi{(\roman{enumi})}
\setcounter{enumi}{3}
\tightlist
\item
  理想 J 是完全割的。
\end{enumerate}

首先注意 (iii) 蕴含 (i)：事实上，在 (iii) 的假设下，代数 \(\mathsf{S}_{\mathrm{A}/\mathrm{J}}(\mathrm{J}/\mathrm{J}^{2})\) 赋予与其分次相伴的拓扑后，在 A/J 上形式光滑 (第 3 节, 例 1)，从而在 k 上形式光滑 (第 2 节, 命题 3, a))；于是断言 (i) 由第 2 节的命题 3, c) 得出。

设 \(\widehat{\mathbf{A}}\) 为代数 A 的分离完备化，\(\widehat{\mathbf{J}}\) 为 J 的分离完备化。典范同态 \(i: \mathbf{A} \to \widehat{\mathbf{A}}\) 诱导出同构 \(\mathrm{A} / \mathrm{J} \longrightarrow \widehat{\mathrm{A}} / \widehat{\mathrm{J}}\) (III, § 2, 第 12 节, 公式 (21))。设 \(\varphi: \mathrm{A} / \mathrm{J} \to \widehat{\mathrm{A}}\) 为该同构的一个提升 (第 4 节, 命题 5)。以 \(\lambda: \widehat{\mathrm{J}} \to \mathrm{J} / \mathrm{J}^2\) 表示由典范同构 \(\mathrm{J} / \mathrm{J}^2 \longrightarrow \widehat{\mathrm{J}} / \widehat{\mathrm{J}}^2\) (III, § 2, 第 12 节, 公式 (21)) 导出的满射。设 \(a\) 为 A 的元素，\(\bar{a}\) 为它在 \(\mathrm{A} / \mathrm{J}\) 中的类，\(z\) 为 \(\widehat{\mathrm{J}}\) 的元素；有 \(\varphi(\bar{a}) \equiv i(a)\) (mod. \(\widehat{\mathrm{J}}\))，从而 \(\varphi(\bar{a})z \equiv i(a)z\) (mod. \(\widehat{\mathrm{J}}^2\))，且 \(\lambda(\varphi(\bar{a})z) = \lambda(i(a)z) = \bar{a}\lambda(z)\)。换言之，\(\lambda\) 是 A/J-线性的（当赋予 \(\widehat{\mathbf{J}}\) 由 \(\varphi\) 诱导的 A/J-模结构时）。

设同态 \(\lambda\) 有一个 A/J-线性截面 \(\sigma : J/J^{2} \longrightarrow \widehat{J}\)。设 \(S\) 为 \(k\)-分次代数 \(S_{A/J}(J/J^{2})\)，\(\widehat{S}\) 为其完备化。设

\[
\theta : \mathsf {S} \to \widehat {\mathsf {A}}
\]

为如下定义的 \(k\)-代数同态：满足 \(\theta(x) = \varphi(x)\)，其中 \(x\) 属于 \(\mathbf{S}^0 = \mathrm{A}/\mathrm{J}\)；且 \(\theta(x) = \sigma(x)\)，其中 \(x\) 属于 \(\mathbf{S}^1 = \mathrm{J}/\mathrm{J}^2\)。由于 \(\theta\) 把 \(\mathbf{S}^1\) 映入 \(\widehat{\mathbf{J}}\)，它把 \(\mathbf{S}^n\) 映入 \(\widehat{\mathbf{J}}^n\)，从而扩张为连续同态 \(\widehat{\boldsymbol{\theta}}: \widehat{\mathbf{S}} \to \widehat{\mathbf{A}}\)。映射 \(\mathrm{gr}_1(\boldsymbol{\theta}): \mathrm{J}/\mathrm{J}^2 \longrightarrow \widehat{\mathrm{J}}/\widehat{\mathrm{J}}^2\) 是 \(\sigma\) 与典范满射 \(\widehat{\mathrm{J}} \longrightarrow \widehat{\mathrm{J}}/\widehat{\mathrm{J}}^2\) 的复合；由于 \(\sigma\) 是 \(\lambda\) 的一个截面，\(\mathrm{gr}_1(\boldsymbol{\theta})\) 与从 \(\mathrm{J}/\mathrm{J}^2\) 到 \(\widehat{\mathrm{J}}/\widehat{\mathrm{J}}^2\) 上的典范同构重合。于是 \(\mathrm{gr}(\boldsymbol{\theta}): \mathbf{S} \to \mathrm{gr}_{\widehat{\mathbf{J}}}(\widehat{\mathbf{A}})\) 是典范满射 \(\beta\) 与典范同构 \(\mathrm{gr}_\mathrm{J}(\mathrm{A}) \to \mathrm{gr}_{\widehat{\mathbf{J}}}(\widehat{\mathbf{A}})\) 的复合 (III, § 2, 第 12 节, 公式 (22))。

现在证明蕴涵 (ii)⇒(iii)。在假设 (ii) 之下，A/J-模 J/J² 是投射的，故 λ 有一个 A/J-线性截面；由前面的构造与该截面相伴的同态 \(\widehat{\boldsymbol{\theta}}:\widehat{\mathbf{S}}\to\widehat{\mathbf{A}}\) 依假设在相伴分次模上诱导出同构，从而是双射 (III, § 2, 第 8 节, 定理 1 的推论 3)，这就证明了 (iii)。

我们来证明 (i)⇒(ii)。设拓扑 k-代数 A 形式光滑。首先证明 A/J-模 \(J/J^{2}\) 是投射的。设 M 为 \(A/J\)-模，\(f:M\to J/J^{2}\) 为满的 A/J-线性映射；目标是证明 f 有一个 A/J-线性截面。

设 \(\pi: A/J^{2} \to A/J\) 为典范满射。由第 2 节的注 1，存在 \(k\)-代数同构 \(\psi: A/J \oplus J/J^{2} \longrightarrow A/J^{2}\)，满足 \(\pi(\psi(y,z)) = y\) 且 \(\psi(0,z) = z\)（对 \(y \in A/J\)、\(z \in J/J^{2}\)）。考察 \(k\)-代数 (A/J) \(\oplus M\) (第 1 节, 例) 与映射 \(u: (A/J) \oplus M \to A/J^{2}\)，它满足 \(u(x,m) = \psi(x,f(m))\)。这是 \(k\)-代数的一个满同态，其核是 \(f\) 在 M 中的核，故是平方为零的幂零理想。典范满射 \(\rho: A \to A/J^{2}\) 连续；由于 A 作为拓扑 \(k\)-代数形式光滑，存在 \(k\)-代数同态 \(\tilde{\rho}: A \to (A/J) \oplus M\) 使 \(u \circ \tilde{\rho} = \rho\)。由于 \(\mathrm{pr}_{1} = \pi \circ \psi = \pi \circ u\)，有 \(\mathrm{pr}_{1} \circ \tilde{\rho} = \pi \circ u \circ \tilde{\rho} = \pi \circ \rho\)，于是 \(\mathrm{pr}_{1} \circ \tilde{\rho}\) 是 A 到 A/J 上的典范满射。因此有 \(\tilde{\rho}(J) \subset M\)，从而 \(\tilde{\rho}(J^{2}) = 0\)，于是 \(\tilde{\rho}\) 诱导出 A/J-线性映射 \(s: J/J^{2} \to M\)。有 \(u \circ \tilde{\rho} = \rho\)，且 \(\mathrm{pr}_{2} \circ \psi^{-1} \circ u(y,m) = f(m)\)（对 \(y \in A/J\) 与 \(m \in M\)）。设 \(x \in J\)，\(\bar{x}\) 为它在 \(J/J^{2}\) 中的类；有 \(f(s(\bar{x})) = f(\mathrm{pr}_{2}(\tilde{\rho}(x))) = \mathrm{pr}_{2}(\psi^{-1}(\bar{x})) = \bar{x}\)。于是 \(s\) 是 \(f\) 的一个截面。

尚需证明同态 \(\beta\) 是单射。由于 A/J-模 \(J/J^{2}\) 是投射的，\(\lambda\) 有一个 A/J-线性截面；以 \(\theta: S \to \widehat{A}\) 表示相伴的同态。同态 gr(θ) 与 \(\beta\) 恒同。设 \(m\) 为整数；以 \(\Sigma_{m}\) 表示分次 \(k\)-代数 \(S\) 模理想 \(\sum_{i > m} S^{i}\) 所得的商，以 \(\theta_{m}: \Sigma_{m} \to A/J^{m+1}\) 表示由 \(\theta\) 导出的同态。\(\theta_{m}\) 与典范满射 \(A/J^{m+1} \longrightarrow A/J\) 的复合是 \(\Sigma_{m}\) 到 \(S^{0} = A/J\) 上的典范投影；于是 \(\theta_{m}\) 的核是幂零理想。由第 4 节的例，存在一个提升 \(\psi_{m}: A \to \Sigma_{m}\)（典范满射 \(A \to A/J^{m+1}\) 的提升）。由于 \(\psi_{m}\) 与从 \(\Sigma_{m}\) 到 A/J 上的典范投影的复合是典范满射，\(\psi_{m}(J)\) 由次数 \(>0\) 的元素组成。通过过渡到相伴分次对象，我们从 \(\psi_{m}\) 得到一个 \(k\)-线性分次映射 gr(\(\psi_{m}\)): gr \(_{J}(A) \to \Sigma_{m}\)，满足 gr \(_{m}(\theta) \circ\) gr \(_{m}(\psi_{m}) = \text{Id}_{J^{m}/J^{m+1}}\)。由此推出 gr \(_{m}(\theta)\)——从而 \(\beta_{m}\) 也——是单射，这就完成了 (ii) 的证明。

最后，当 A 是 Noether 的时，条件 (ii) 与 (iv) 等价 (§ 5, 第 2 节, 定理 1)。

推论 1.--- 设 \(k\) 为域，A 为 \(k\)-局部 Noether 代数，且扩张 \(\kappa_{A}\)（在 \(k\) 上）可分。下列条件等价：

\begin{enumerate}
\def\labelenumi{(\roman{enumi})}
\item
  \(k\)-代数 A 对 \(\mathfrak{m}_{\mathrm{A}}\)-进拓扑形式光滑；
\item
  环 A 正则；
\item
  \(k\)-代数 A 绝对正则 (\(\S 6, \mathrm{n}^{\circ} 4\), 定义 1)；
\item
  k-代数 \(\hat{A}\) 同构于 \(\kappa_{A}[[T_{1},\ldots,T_{n}]]\)，其中 \(n=\dim A\)。
\end{enumerate}

条件 (ii) 与 (iii) 的等价由 § 6 第 4 节的例 3 得出，且两者等价于说理想 \(\mathfrak{m}_{\mathrm{A}}\) 是完全割的 (VIII, § 5, 第 2 节, 定理 1)。此外，\(\widehat{\mathrm{A}}\) 到 \(\kappa_{\mathrm{A}}[[\mathrm{T}_{1},\ldots,\mathrm{T}_{n}]]\) 上的每个环同构都是双连续的，因为它们都是局部环。由于 \(k\)-代数 \(\mathrm{A}/\mathfrak{m}_{\mathrm{A}}\) 形式光滑 (第 3 节, 定理 1)，取 \(\mathrm{J}=\mathfrak{m}_{\mathrm{A}}\) 应用定理 2 即得推论。

推论 2.--- 设 k 为域，A 为 Noether k-代数，J 为 A 的含于 A 的根式中的理想。设 k-代数 A 对 J-进拓扑形式光滑。则它绝对正则。

事实上，设 \(k'\) 为 \(k\) 的有限扩张，\(A'\) 为 A-代数 \(A_{(k')}\)；尚需证明：对每个极大理想 \(\mathfrak{m}'\)（\(A'\) 的），Noether 局部环 \(A_{\mathfrak{m}'}'\) 都是正则的。我们有 \(JA' \subset \mathfrak{m}'\)：事实上，\(\mathfrak{m}'\) 在 A 中的逆像是 A 的一个极大理想 (V, § 2, 第 1 节, 命题 1)，故包含 J。\(k'\)-代数 \(A'\) 对 \(JA'\)-进拓扑形式光滑 (第 2 节, 命题 4, b))，而 \(k'\)-代数 \(A_{\mathfrak{m}'}'\) 对 \(JA_{\mathfrak{m}'}'\)-进拓扑形式光滑 (第 2 节, 命题 4, a))，故对 \(\mathfrak{m}'A_{\mathfrak{m}'}'\)-进拓扑也形式光滑。设 \(k_0\) 为 \(k'\) 的素子域。则 \(A_{\mathfrak{m}'}'\) 在 \(k_0\) 上对 \(\mathfrak{m}'A_{\mathfrak{m}'}'\)-进拓扑形式光滑 (第 3 节定理 1 的推论)；由于 \(\kappa(\mathfrak{m}')\) 在 \(k_0\) 上可分，环 \(A_{\mathfrak{m}'}'\) 正则 (推论 1)。

推论 3.--- 设 k 为环，A 为形式光滑 k-代数。

\begin{enumerate}
\def\labelenumi{\alph{enumi})}
\item
  A-模 \(\Omega_k(A)\) 是投射的。
\item
  设环 A ⊗k A 是 Noether 的。设 μ : A ⊗k A → A 为满足 μ(x ⊗ y) = xy 的同态；则理想 Ker(μ) 是完全割的。
\end{enumerate}

\(k\)-代数 A 与 \(\mathbf{A} \otimes_{k} \mathbf{A}\) 形式光滑 (第 2 节, 命题 4, c))，且 A 同构于 \(\mathbf{A} \otimes_{k} \mathbf{A}\) 模 \(\mu\) 的核 I 所得的商。依定义，\(\Omega_{k}(\mathbf{A}) = \mathbf{I}/\mathbf{I}^{2}\)。于是 a) 与 b) 由定理 2 得出。

